% 13-declaraciones.tex — Fuente: docs/informe-final/13-declaraciones-referencias.md
% (sección "Declaraciones obligatorias" únicamente; la bibliografía se
% generó a partir de referencias.bib y se imprime automáticamente en
% main.tex con \bibliography{referencias}, ver también el Capítulo de
% Anexos para las tablas de clasificación de referencias.)

\chapter{Declaraciones obligatorias}
\label{cap:declaraciones}

\section{Disponibilidad de código}

Repositorio público:
\url{https://github.com/JohanCarvajal04/Proyecto-WEB-ARTISYNC}. Etiqueta
de esta entrega: \texttt{v1.0.0}. Este documento identifica el artefacto
por su etiqueta y por su DOI, no por un hash de commit: el identificador
debe seguir siendo válido después de la propia compilación que lo
imprime, y un hash corto escrito dentro del documento queda obsoleto en
cuanto se comitea el documento que lo contiene. Para obtener el commit
exacto que la etiqueta designa en cualquier momento, basta
\texttt{git rev-list -n 1 v1.0.0} sobre el repositorio público; la
etiqueta local y la de \texttt{origin} están verificadas como
coincidentes.

\paragraph{Relación entre la etiqueta y el depósito de Zenodo.} El
depósito Zenodo del software (\texttt{10.5281/zenodo.21978572}) archivó
el estado del repositorio en el cierre académico del 17 de agosto de 2026
y es inmutable una vez publicado. La etiqueta \texttt{v1.0.0} se
reposicionó con posterioridad para incorporar el trabajo de corrección de
observaciones realizado hasta el cierre del examen final, de modo que el
\emph{snapshot} de Zenodo y el commit que hoy designa la etiqueta
\textbf{no son el mismo commit}: el primero es un archivo histórico
citable, la segunda señala el artefacto evaluado. Ambos son alcanzables
desde el repositorio público y la diferencia entre ellos es exactamente
el historial de correcciones documentado en \texttt{CHANGELOG.md} y en
\texttt{docs/observaciones/OBSERVACIONES.md}.

El tag \texttt{v1.0.0} de la copia local del repositorio llegó a apuntar,
por un tiempo, a un commit distinto (\texttt{5b61f86}), creado por
separado; se verificó y corrigió el 01-09-2026 realineando el tag local
con el de GitHub (\texttt{git fetch origin tag v1.0.0}), sin alterar
ningún artefacto ya publicado. El trabajo de consolidación posterior al
cierre académico de la Entrega Final (refactor de permisos y
endurecimiento de seguridad) se etiquetó aparte como \texttt{v1.1.0},
fuera del alcance de esta entrega. Ver
\texttt{docs/observaciones/OBSERVACIONES.md}, OBS-AUTO-12.
DOI Zenodo del software:
\texttt{10.5281/zenodo.21978572} (archivo de \texttt{v1.0.0}; la versión
anterior, \texttt{v0.9.0-rc}, quedó archivada aparte en
\texttt{10.5281/zenodo.21730559}).

\section{Disponibilidad de datos}

Depósito Zenodo separado del software para el dataset de
mediciones (\texttt{docs/mediciones/}): DOI
\texttt{10.5281/zenodo.22236251}, licencia Creative Commons
Attribution 4.0 International (CC~BY~4.0), siguiendo el principio de
citación independiente de software y datos. Contiene los datos crudos
de rendimiento (k6), seguridad (OWASP, ZAP, análisis estático),
usabilidad (SUS), cobertura (JaCoCo) y calidad web (Lighthouse), junto
con \texttt{DATA-DICTIONARY.md} y \texttt{DATA-PROVENANCE.md}. Los
mismos datos crudos siguen disponibles en el repositorio bajo
\texttt{docs/mediciones/}, referenciados en la matriz de trazabilidad.

\section{Disponibilidad de materiales}

\begin{itemize}
  \item Colección Postman (26 peticiones):
    \texttt{Pruebas.postman\_collection.json}.
  \item Script de carga k6: \texttt{k6/catalogo-load.js}.
  \item Configuración Lighthouse:
    \texttt{artisync/Frontend/lighthouserc.mobile.json},
    \texttt{lighthouserc.desktop.json}.
  \item Scripts de análisis SUS: \texttt{docs/mediciones/sus/analisis-sus.py},
    \texttt{graficar-sus.py}.
  \item Protocolo SUS y plantilla de consentimiento:
    \texttt{docs/mediciones/sus/instrucciones-formulario.md},
    \texttt{docs/etica/consentimientos/plantilla.md}.
\end{itemize}

\section{Declaración ética}

Ver \texttt{docs/etica/ETHICS.md} para el detalle completo: resguardo de
consentimientos informados fuera del repositorio público, participación
voluntaria sin compensación, anonimización por código de participante
(\texttt{P01}--\texttt{P16}), sin grabación de audio/video.

\section{Contribuciones de autores (taxonomía CRediT)}

Ver la matriz completa de los catorce roles CRediT en
\texttt{CONTRIBUTORS.md}, con asignación definitiva para v1.0.0. Bone
Arroyo, Niurca Scarleth, colabora con el equipo por ser compañera de
curso en la materia de Administración de Bases de Datos, en la que este
mismo proyecto también se utiliza como caso de estudio; su aporte se
concentra en la capa de base de datos (procedimientos almacenados y
funciones SQL) y tareas de backend relacionadas. Esta colaboración
cruzada fue comunicada al docente-director durante las sesiones de
clase. Resumen:
los cuatro integrantes (Bone Arroyo N., Carvajal Loor J., Figueroa
Morales B., Rios Cuyabazo J.) comparten Conceptualización, Metodología,
Software, Validación, Investigación, Recursos, Redacción (borrador y
revisión) y Administración del proyecto; Análisis formal se concentra en
Figueroa B. y Rios J.; Curación de datos en Rios J.; Visualización en
Bone N. y Carvajal J.; Supervisión corresponde al docente-director del
PFC, Dr. Gleiston Cicerón Guerrero Ulloa; y Adquisición de fondos se
declara explícitamente sin asignar, pues el proyecto se desarrolló con
recursos propios del equipo y niveles gratuitos/académicos de
proveedores cloud, sin financiamiento externo que gestionar.

\section{Declaración de conflictos de interés}

Ver \texttt{docs/etica/ETHICS.md}: el equipo declara ausencia de
conflictos de interés.

\section{Declaración de financiamiento}

Ver \texttt{docs/etica/ETHICS.md}: sin financiamiento externo;
desarrollado con recursos propios del equipo y niveles gratuitos/académicos
de proveedores cloud.

\section{Declaración de uso de asistentes de inteligencia artificial}

Ver \texttt{docs/etica/ai-disclosure.md}. Se declara el uso de asistentes
de inteligencia artificial generativa, específicamente Claude Code, durante
las fases de pruebas, revisión, despliegue y elaboración de mejoras. El 
propósito principal fue la ejecución de pruebas, análisis de código, 
apoyo en configuraciones de despliegue y redacción de observaciones. 
Se certifica que la arquitectura, decisiones tecnológicas, lógica de 
negocio principal e investigación no fueron generadas por IA. Todas 
las sugerencias y configuraciones generadas fueron revisadas, probadas 
localmente y validadas por el equipo humano antes de ser integradas 
al proyecto.

\section{Agradecimientos}

Agradecemos sinceramente a todas las personas e instituciones que han hecho
posible el desarrollo de este proyecto. Esta experiencia ha sido invaluable para
nuestro crecimiento tanto profesional como personal, brindándonos la oportunidad
de profundizar nuestros conocimientos prácticos y aprender más sobre el
desarrollo de software, las arquitecturas escalables y el diseño centrado en el
usuario.

Queremos extender un profundo agradecimiento a todos los usuarios que, de forma
anónima y desinteresada, participaron en nuestras pruebas de usabilidad
(evaluación SUS). Su tiempo, disposición y valiosa retroalimentación colectiva
fueron fundamentales para iterar y mejorar significativamente la experiencia de
la plataforma Artisync.

Finalmente, agradecemos a nuestros docentes y a la institución por la formación,
guía y apoyo constante a lo largo de todo este proceso de aprendizaje.
